\documentclass[a4paper,11pt]{jsarticle}
\usepackage[dvipdfmx]{graphicx}
\usepackage{geometry}
\usepackage{amsmath}
\usepackage{booktabs}
\usepackage{longtable}
\usepackage{array}
\usepackage{siunitx}
\usepackage{url}
\geometry{margin=20mm}

\title{p2MEG 最終解析における ACC constraint 付き尤度解析\\
step4f\_run1to20\_eg\_doubleonly\_allpatterns\_500ns.txt}
\author{p2MEG 解析レポート}
\date{2026年3月8日}

\begin{document}
\maketitle

\begin{abstract}
最終 5D データ
\texttt{data/finaldata/step4f\_run1to20\_eg\_doubleonly\_allpatterns\_500ns.txt}
に対し、TSB から AW 内 ACC 数を見積もり、その結果を Gaussian constraint として
尤度関数に導入した。

今回の重要な更新点は、ACC PDF の時間項 \(p_t(t|E_\gamma)\) を
TSB 事象の raw ヒストグラムから直接作るのではなく、
TSB のみを使って fit した対称関数
\[
f(t)=A\exp\left(-\frac{t^2}{2\sigma^2}\right)+C
\]
を AW 内へ外挿して構成したことである。
この fit は ACC 数見積もりに使ったものと
\textbf{同じ仕様、同じ実装}を使い回している。
これにより、AW 内事象を time template に使わずに、
AW 内でも有限の時間 PDF を持つ ACC モデルを作ることができた。

AW 内 ACC 数の見積もりは
\[
N_{\mathrm{ACC,pred}} = 10.13 \pm 8.75
\]
であり、これを Gaussian constraint として NLL に加えた。
その結果、AW 内 18 事象に対し
\begin{align}
N_{\mathrm{sig}} &= 3.93\times 10^{-9},\\
N_{\mathrm{rmd}} &= 2.71 \pm 2.56,\\
N_{\mathrm{acc}} &= 14.47 \pm 4.93
\end{align}
を得た。
現状では signal 成分は不要であり、観測事象は主として ACC と少量の RMD で説明される。
\end{abstract}

\section{目的}
本レポートの目的は、次の一連の処理を 1 つにまとめて記録することである。
\begin{enumerate}
  \item TSB から AW 内 ACC 数を見積もる。
  \item その結果を constraint として尤度関数へ入れる。
  \item ACC PDF の時間項を、ACC 数見積もりと同じ fit 仕様で作る。
  \item 変更後の ACC PDF を使って最尤解析を行う。
  \item 手順、式、図、生成した PDF の作り方まで含めてまとめる。
\end{enumerate}

\section{入力データと解析窓}
入力は
\[
\texttt{data/finaldata/step4f\_run1to20\_eg\_doubleonly\_allpatterns\_500ns.txt}
\]
である。
各行は
\[
(E_e,\ E_\gamma,\ t,\ \phi_e,\ \phi_\gamma)
\]
の 5 変数を持つ。

解析窓は
\begin{align}
20 &\le E_e \le 80 \ \mathrm{MeV},\\
20 &\le E_\gamma \le 80 \ \mathrm{MeV},\\
-50 &\le t \le 50 \ \mathrm{ns},\\
1.7 &\le \theta_{eg} \le \pi
\end{align}
とした。
ここで
\[
\theta_{eg}=|\phi_e-\phi_\gamma|
\]
である。

また TSB は
\[
t\in[-500,500]\ \mathrm{ns}
\quad\text{かつ}\quad
t\notin[-50,50]\ \mathrm{ns}
\]
と定義した。

\section{AW 内 ACC 数の見積もり}
\subsection{見積もりの方針}
AW 内 ACC 数の見積もりには
\texttt{macros/fit\_acc\_time\_by\_eg.C}
と
\texttt{macros/predict\_acc\_from\_tsb\_by\_eg.C}
を用いた。

ここでは \(E_\gamma\) を
\[
[0,10],\ [10,20],\ [20,30],\ [30,80]\ \mathrm{MeV}
\]
の 4 bin に分け、
各 bin ごとに時間分布を
\[
f_i(t)=A_i\exp\left(-\frac{t^2}{2\sigma_i^2}\right)+C_i
\]
で fit した。

ただし fit に使う時間ヒストグラムには、4D 解析窓
\[
(E_e,E_\gamma,\theta_{eg},t)\in\mathrm{AW}
\]
の事象を入れない。
したがって、時間形状は AW 外のデータだけで学習される。

ここで使う fit の実装は
\texttt{include/p2meg/AccTimeFit.h}
に共通化した。
\texttt{macros/fit\_acc\_time\_by\_eg.C} でも
\texttt{src/MakeACCGridPdf.cc} でも、同じ
\texttt{AccTimeFit\_Run(...)} を呼ぶ。
したがって
\begin{itemize}
  \item model は \(A\exp(-t^2/2\sigma^2)+C\)
  \item reject window は \(|t|<20\) ns
  \item \(\sigma\) の下限は 20 ns
\end{itemize}
が両者で完全に一致している。

\subsection{fit と外挿}
\(|t|<20\) ns は blind core とし、
\[
\sigma_i \ge 20\ \mathrm{ns}
\]
を課した。
fit 後、
\begin{equation}
R_i^{\mathrm{fit}}=
\frac{\int_{-50}^{50} f_i(t)\,dt}
{\int_{-500}^{-50} f_i(t)\,dt+\int_{50}^{500} f_i(t)\,dt}
\end{equation}
を定義した。

また、生カウントから
\[
R_i^{\mathrm{data}} = \frac{N_i(\mathrm{AW})}{N_i(\mathrm{TSB})}
\]
も作った。

最後に
\begin{equation}
N_{\mathrm{ACC,pred},i} = R_i^{\mathrm{fit}}N_i(\mathrm{TSB})
\end{equation}
で bin ごとの AW 内 ACC 数を予測し、全 bin を足して
\[
N_{\mathrm{ACC,pred}}(\mathrm{AW})
\]
を得た。

\subsection{数値結果}
今回の結果は
\begin{align}
N_{\mathrm{AW,obs}} &= 18,\\
N_{\mathrm{ACC,pred}} &= 10.1275 \pm 8.7521,\\
N_{\mathrm{AW,obs}} - N_{\mathrm{ACC,pred}} &= 7.8725
\end{align}
である。

主要な bin の内訳は
\begin{center}
\begin{tabular}{ccccc}
\toprule
\(E_\gamma\) [MeV] & \(N_{\mathrm{AW}}\) & \(N_{\mathrm{TSB}}\) & \(R_i^{\mathrm{fit}}\) & \(N_{\mathrm{ACC,pred},i}\) \\
\midrule
20--30 & 16 & 63 & 0.116662 & 7.3497 \\
30--80 &  2 & 25 & 0.111111 & 2.7778 \\
\bottomrule
\end{tabular}
\end{center}
であった。

\begin{figure}[p]
  \centering
  \includegraphics[width=0.95\linewidth,page=1]{predict_acc_from_tsb_by_eg_step4f_run1to20_eg_doubleonly_allpatterns_500ns.pdf}
  \vspace{3mm}
  \includegraphics[width=0.95\linewidth,page=2]{predict_acc_from_tsb_by_eg_step4f_run1to20_eg_doubleonly_allpatterns_500ns.pdf}
  \caption{TSB から AW 内 ACC 数を予測した結果。}
  \label{fig:predacc}
\end{figure}

\begin{figure}[p]
  \centering
  \includegraphics[width=0.95\linewidth,page=1]{fit_acc_time_by_eg_step4f_run1to20_eg_doubleonly_allpatterns_500ns.pdf}
  \caption{ACC 時間 fit の meta page。入力、選別、blind core、4D AW 除外の条件をまとめた。}
  \label{fig:fitaccmeta}
\end{figure}

\begin{figure}[p]
  \centering
  \includegraphics[width=0.95\linewidth,page=2]{fit_acc_time_by_eg_step4f_run1to20_eg_doubleonly_allpatterns_500ns.pdf}
  \vspace{2mm}
  \includegraphics[width=0.95\linewidth,page=3]{fit_acc_time_by_eg_step4f_run1to20_eg_doubleonly_allpatterns_500ns.pdf}
  \caption{\(E_\gamma=[0,10]\), \([10,20]\) MeV の時間 fit。}
  \label{fig:fitacc12}
\end{figure}

\begin{figure}[p]
  \centering
  \includegraphics[width=0.95\linewidth,page=4]{fit_acc_time_by_eg_step4f_run1to20_eg_doubleonly_allpatterns_500ns.pdf}
  \vspace{2mm}
  \includegraphics[width=0.95\linewidth,page=5]{fit_acc_time_by_eg_step4f_run1to20_eg_doubleonly_allpatterns_500ns.pdf}
  \caption{\(E_\gamma=[20,30]\), \([30,80]\) MeV の時間 fit。}
  \label{fig:fitacc34}
\end{figure}

\section{ACC PDF の作り方}
\subsection{4D 形状 \(p_4(E_e,E_\gamma,\phi_e,\phi_\gamma)\)}
ACC PDF の 4D 本体は
\texttt{src/MakeACCGridPdf.cc}
で作った。
手順は次の通りである。
\begin{enumerate}
  \item \((E_e,E_\gamma,\theta_{eg})\) が解析窓内の事象を選ぶ。
  \item そのうち \(t\) が TSB にある事象だけを採用する。
  \item \(h_E(E_e,\phi_e)\) と \(h_G(E_\gamma,\phi_\gamma)\) を作る。
  \item \(E\) 軸方向のみ平滑化し、密度に正規化する。
  \item
  \[
  p_4(E_e,E_\gamma,\phi_e,\phi_\gamma)
  \propto
  p_E(E_e,\phi_e)\,p_G(E_\gamma,\phi_\gamma)
  \]
  として因子化した 4D 格子 PDF を作る。
\end{enumerate}

今回の生成では
\begin{align}
N_{\mathrm{events,total}} &= 1578,\\
N_{\mathrm{in\ 3D\ window}} &= 106,\\
N_{\mathrm{TSB}} &= 88
\end{align}
であり、4D PDF の形状情報は TSB の 88 事象のみから作られている。

\subsection{時間項 \(p_t(t|E_\gamma)\) を fit 関数から作る}
今回の更新点はここである。
以前は time template \texttt{tshape} を raw ヒストグラムで保存していたため、
4D AW 内事象を template からも除外すると、
AW 内の bin content が 0 になり、
\[
\int_{\mathrm{AW}} \mathrm{tshape}(t)\,dt = 0
\]
となっていた。
その結果、\texttt{ACCGridPdf\_Load} では時間項を正規化できず、
AW 内一様分布へフォールバックしていた。

今回はこの問題を避けるため、\texttt{MakeACCGridPdf.cc} の中で
time template を次の手順で作るように変更した。
\begin{enumerate}
  \item raw の TSB-only 時間ヒストグラム \(h_t^{\mathrm{raw}}\) を作る。
  \item 4D AW 内事象は \(h_t^{\mathrm{raw}}\) に入れない。
  \item \(|t|<20\) ns の blind core を \textbf{fit から reject} する。
  \item TSB 領域のみを用いて
  \[
  f(t)=A\exp\left(-\frac{t^2}{2\sigma^2}\right)+C
  \]
  を fit する。
  \item 得られた \(f(t)\) を用いて、全時間域 \([-500,500]\) の各 bin に
  \[
  \frac{1}{\Delta t_{\mathrm{bin}}}
  \int_{\mathrm{bin}} f(t)\,dt
  \]
  を入れ、fit-based template ヒストグラムを作る。
\end{enumerate}

ここで重要なのは、これは ACC 数見積もりで使った fit と
\textbf{全く同じ仕様}だということである。
すなわち、PDF 用だけ別の時間 model を導入したのではなく、
ACC 数予測で採用した時間 model をそのまま
\(p_t(t|E_\gamma)\) に再利用している。

この方法では、AW 内 bin にデータを直接入れなくても、
TSB fit の外挿により AW 内で有限の密度が得られる。
したがって
\[
\int_{\mathrm{AW}} \mathrm{tshape}_{\mathrm{fit}}(t)\,dt > 0
\]
となり、\texttt{ACCGridPdf\_Load} は
AW 内正規化済みの時間 PDF
\[
p_t(t|E_\gamma)=
\frac{\mathrm{tshape}_{\mathrm{fit}}(t)}
{\int_{\mathrm{AW}}\mathrm{tshape}_{\mathrm{fit}}(t)\,dt}
\]
を使えるようになる。

ここで重要なのは、
\begin{quote}
AW 内で有限の密度を持つが、その値は AW データからではなく
TSB fit の外挿から得ている
\end{quote}
という点である。

\subsection{今回使った最終 ACC PDF}
以上より、今回の尤度 fit で使った ACC PDF は
\begin{equation}
p_{\mathrm{acc}}(x)
=
p_4(E_e,E_\gamma,\phi_e,\phi_\gamma)\times p_t(t|E_\gamma)
\end{equation}
であり、
\begin{itemize}
  \item 4D 形状 \(p_4\) は TSB 事象から構成
  \item 時間項 \(p_t\) は TSB-only fit 関数から構成
\end{itemize}
という意味で、AW 内事象を直接学習に使っていない。

\begin{figure}[p]
  \centering
  \includegraphics[width=0.95\linewidth,page=1]{acc_grid_hist_step4f_doubleonly_fitbased.pdf}
  \caption{fit-based time template を持つ ACC PDF の meta page。}
  \label{fig:accgridmeta}
\end{figure}

\begin{figure}[p]
  \centering
  \includegraphics[width=0.95\linewidth,page=2]{acc_grid_hist_step4f_doubleonly_fitbased.pdf}
  \vspace{2mm}
  \includegraphics[width=0.95\linewidth,page=3]{acc_grid_hist_step4f_doubleonly_fitbased.pdf}
  \caption{ACC PDF の 1D/2D 可視化。\(t\) 分布も一様ではなく fit-based template 由来の形を持つ。}
  \label{fig:accgridshape}
\end{figure}

\section{尤度関数と constraint}
\subsection{NLL}
尤度 fit は
\texttt{scripts/run\_nll\_fit.cc}
で行った。
解析窓内 18 事象だけを使い、signal, RMD, ACC の 3 成分で fit した。

NLL は
\begin{equation}
\mathrm{NLL}
=
\left(\sum_k N_k\right)
-\sum_i \log\left(\sum_k N_k p_k(x_i)\right)
+\mathrm{ConstraintNLL}(N_k)
\end{equation}
である。

\subsection{ACC constraint}
\texttt{src/ConstraintNLL.cc}
に
\[
N_{\mathrm{acc}} = 10.1275 \pm 8.75205
\]
の Gaussian constraint を入れた。
したがって ACC 制約項は
\begin{equation}
\mathrm{Constraint}_{\mathrm{acc}}
=
\frac{1}{2}
\left(
\frac{N_{\mathrm{acc}}-10.1275}{8.75205}
\right)^2
\end{equation}
である。

signal と RMD には従来通り広い constraint を残し、
実質的には自由に動けるようにした。

\section{制約付き最尤解析の結果}
\subsection{結果}
新しい fit-based ACC 時間項を使って
\texttt{run\_nll\_fit} を実行した結果、
\begin{align}
\mathrm{status} &= 0,\\
\mathrm{NLL}_{\min} &= 193.71,\\
N_{\mathrm{sig}} &= 3.92803\times 10^{-9},\\
N_{\mathrm{rmd}} &= 2.71048 \pm 2.56384,\\
N_{\mathrm{acc}} &= 14.4692 \pm 4.92959
\end{align}
を得た。

\begin{table}[htbp]
\centering
\caption{ACC constraint 付き尤度 fit の結果}
\label{tab:nllresult}
\begin{tabular}{lc}
\toprule
量 & 値 \\
\midrule
解析窓内事象数 & 18 \\
fit status & 0 \\
\(\mathrm{NLL}_{\min}\) & 193.71 \\
\(N_{\mathrm{sig}}\) & \(3.92803\times 10^{-9}\) \\
\(N_{\mathrm{rmd}}\) & \(2.71048 \pm 2.56384\) \\
\(N_{\mathrm{acc}}\) & \(14.4692 \pm 4.92959\) \\
\bottomrule
\end{tabular}
\end{table}

\subsection{解釈}
constraint の中心値 \(10.13\) 事象に対し、fit は
\[
N_{\mathrm{acc}}\simeq 14.47
\]
を選んだ。
これは \(\pm 8.75\) の constraint 幅の中では十分許容される。
残りは
\[
N_{\mathrm{rmd}}\simeq 2.71
\]
で補われ、signal は不要である。

したがって、fit-based な ACC 時間項を使っても結論は変わらず、
\begin{quote}
現時点の AW 18 事象は、主として ACC と少量の RMD で説明できる
\end{quote}
と判断される。

\section{実行手順}
\subsection{ACC 数の見積もり}
\begin{verbatim}
root -l -b -q 'macros/fit_acc_time_by_eg.C(
  "data/finaldata/step4f_run1to20_eg_doubleonly_allpatterns_500ns.txt")'

root -l -b -q 'macros/predict_acc_from_tsb_by_eg.C(
  "data/finaldata/step4f_run1to20_eg_doubleonly_allpatterns_500ns.txt")'
\end{verbatim}

\subsection{fit-based ACC PDF の生成}
\begin{verbatim}
./build/make_acc_grid_pdf \
  data/finaldata/step4f_run1to20_eg_doubleonly_allpatterns_500ns.txt \
  data/pdf_cache/acc_grid.root acc_grid
\end{verbatim}

\subsection{制約付き最尤解析}
\begin{verbatim}
./build/run_nll_fit \
  data/finaldata/step4f_run1to20_eg_doubleonly_allpatterns_500ns.txt
\end{verbatim}

\subsection{ACC PDF 図の生成}
\begin{verbatim}
./build/plot_acc_grid_pdf \
  data/pdf_cache/acc_grid.root acc_grid \
  doc/finalanalysis/acc_grid_hist_step4f_doubleonly_fitbased.pdf
\end{verbatim}

\section{このレポート PDF の作り方}
このレポート自身は
\texttt{doc/finalanalysis/report\_acc\_constraint\_nll\_step4f.tex}
として保存し、
\begin{verbatim}
cd doc/finalanalysis
platex report_acc_constraint_nll_step4f.tex
platex report_acc_constraint_nll_step4f.tex
dvipdfmx report_acc_constraint_nll_step4f.dvi
\end{verbatim}
で PDF 化した。

図としては
\begin{itemize}
  \item \texttt{fit\_acc\_time\_by\_eg\_*.pdf}
  \item \texttt{predict\_acc\_from\_tsb\_by\_eg\_*.pdf}
  \item \texttt{acc\_grid\_hist\_step4f\_doubleonly\_fitbased.pdf}
\end{itemize}
をそのまま埋め込んでいる。
したがって、このレポート 1 本を開けば
\begin{itemize}
  \item ACC 数見積もりの fit 図
  \item TSB からの AW 予測図
  \item ACC PDF の可視化図
  \item 制約付き尤度解析の結果
\end{itemize}
をまとめて確認できる。

\section{まとめ}
TSB から見積もった
\[
N_{\mathrm{ACC,pred}} = 10.13 \pm 8.75
\]
を Gaussian constraint として導入し、さらに ACC PDF の時間項を
TSB-only fit 関数から構成するように変更した。
このとき、PDF の時間 fit は ACC 数見積もりと同じ仕様・同じ実装を使った。

その結果、
\[
N_{\mathrm{sig}}\simeq 0,\qquad
N_{\mathrm{rmd}}\simeq 2.71,\qquad
N_{\mathrm{acc}}\simeq 14.47
\]
を得た。

したがって、時間項を fit-based に改善した後でも、
現在の最終解析窓の事象は signal を必要とせず、
ACC 優勢の描像と整合している。

\end{document}
